\scanpage{47}
% slanted
\margin{\unsure{remonter|remonte|remarque} en cor. à 3.1.4 ?}
3.1.16. L'isomorphisme \deleted{d'induction}\inserted{d'adjonction} s'écrit
\[
 \mathrm{Hom}(Rf_!K,L)\simeq\mathrm{Hom}(K,Rf^!L)
 \tag*{(3.1.16.1)}
\]
\inserted{avec} \deleted{\{} $K\in\mathrm{Ob}\,D(X,f^*\mathcal A)$ et $L\in\mathrm{Ob}\,D^+(S,\mathcal A)$. \deleted{\}} Pour $K\in\mathrm{Ob}\,D^-(X,f^*\mathcal A)$ et $L\in\mathrm{Ob}\,D^+(S,f^*\mathcal A)$\ensuremath{{}_{*}} \deleted{cet isomorphisme} complexe d'injectifs, cet isomorphisme s'écrit, au niveau des complexes,
\margin{$L$, ou $K$ et $L$ ?}
\[
 \mathrm{Hom}^\bullet(f_!^\bullet K,L)\simeq
 \mathrm{Hom}^\bullet(K,f_\bullet^!L).
\]
Pour tout \deleted{$V\in\mathrm{Ob}\,S_{\mathrm{et}}$}\inserted{$j:V\to S$ dans $S_{\mathrm{et}}$}, on dispose de
\[
\begin{aligned}
 \mathrm{Hom}_{f^*V}^\bullet(K,f_\bullet^!L)
 &=\mathrm{Hom}^\bullet(j'_!j'^*K,f_\bullet^!L)
 \underset{\sim}{\longleftrightarrow}\mathrm{Hom}^\bullet(f_!^\bullet(j'_!j'^*K),L)\longrightarrow\\
 &\deleted{\longrightarrow\mathrm{Hom}\ \mathrm{Hom}}\\
 &\longrightarrow\mathrm{Hom}(j_!j^*(f_!^\bullet K),L)
 =\mathrm{Hom}_{\deleted{\unsure{S}}V}(f_!^\bullet K,L)
\end{aligned}
\]
d'où un morphisme
\[
 Rf_*R\underline{\mathrm{Hom}}(K,Rf^!L)\xrightarrow{\sim}
 R\underline{\mathrm{Hom}}(Rf_!K,L)
 \tag*{(3.1.16.2)}
\]
% margin note written sideways
\margin{devrait remonter en cor. 3.1.4.8 à 3.1.4 (conséquence triviale de 3.1.4.7)}
que par localisation de 3.1.16.1 on vérifie être un isomorphisme.

3.1.17. \uline{Exemple} :
% slanted
\margin{\unsure{Videz|Vider}, ou en faire \unsure{remarque|rem.|venir} d'une ligne.}
Soi\inserted{en}t $X$ un schéma compactifiable sur le spectre d'un corps algébriquement clos $k$ \deleted{\unsure{et}}, $x$ un point fermé de $X$ et $n$ un entier $\geq1$ inversible dans $k$\unsure{,|.} La formule (3.1.11.1) fournit, puisque $\mathbb Z/n\mathbb Z$ est un $\mathbb Z/n\mathbb Z$-module injectif ; \deleted{Pour $U(x)$ parcourant les voisinages étales de $x$}
\[
 \underline H^{-i}(Rf^!\mathbb Z/n\mathbb Z)_x
 \simeq\underset{U(x)}{\underrightarrow{\lim}}\,
 \mathrm{Hom}(\deleted{\unsure{\mathbb Z/}}H_c^i(U,\mathbb Z/n\mathbb Z),\mathbb Z/n\mathbb Z),
\]
\inserted{où $U(x)$ est la catégorie des voisinages étales de $x$ dans $X$.}
% slanted
\margin{pourquoi pas $W(x),W$ ?}
Il résulte vraisemblablement du théorème de bidualité (SGA 5 I \qquad) (qui dépend de la résolution des singularités)
\margin{?}
% archive sheet number: 46
